% Options for packages loaded elsewhere
\PassOptionsToPackage{unicode}{hyperref}
\PassOptionsToPackage{hyphens}{url}
%
\documentclass[
]{article}
\usepackage{amsmath,amssymb}
\usepackage{lmodern}
\usepackage{iftex}
\ifPDFTeX
  \usepackage[T1]{fontenc}
  \usepackage[utf8]{inputenc}
  \usepackage{textcomp} % provide euro and other symbols
\else % if luatex or xetex
  \usepackage{unicode-math}
  \defaultfontfeatures{Scale=MatchLowercase}
  \defaultfontfeatures[\rmfamily]{Ligatures=TeX,Scale=1}
\fi
% Use upquote if available, for straight quotes in verbatim environments
\IfFileExists{upquote.sty}{\usepackage{upquote}}{}
\IfFileExists{microtype.sty}{% use microtype if available
  \usepackage[]{microtype}
  \UseMicrotypeSet[protrusion]{basicmath} % disable protrusion for tt fonts
}{}
\makeatletter
\@ifundefined{KOMAClassName}{% if non-KOMA class
  \IfFileExists{parskip.sty}{%
    \usepackage{parskip}
  }{% else
    \setlength{\parindent}{0pt}
    \setlength{\parskip}{6pt plus 2pt minus 1pt}}
}{% if KOMA class
  \KOMAoptions{parskip=half}}
\makeatother
\usepackage{xcolor}
\usepackage[margin=1in]{geometry}
\usepackage{graphicx}
\makeatletter
\def\maxwidth{\ifdim\Gin@nat@width>\linewidth\linewidth\else\Gin@nat@width\fi}
\def\maxheight{\ifdim\Gin@nat@height>\textheight\textheight\else\Gin@nat@height\fi}
\makeatother
% Scale images if necessary, so that they will not overflow the page
% margins by default, and it is still possible to overwrite the defaults
% using explicit options in \includegraphics[width, height, ...]{}
\setkeys{Gin}{width=\maxwidth,height=\maxheight,keepaspectratio}
% Set default figure placement to htbp
\makeatletter
\def\fps@figure{htbp}
\makeatother
\setlength{\emergencystretch}{3em} % prevent overfull lines
\providecommand{\tightlist}{%
  \setlength{\itemsep}{0pt}\setlength{\parskip}{0pt}}
\setcounter{secnumdepth}{-\maxdimen} % remove section numbering
\ifLuaTeX
  \usepackage{selnolig}  % disable illegal ligatures
\fi
\IfFileExists{bookmark.sty}{\usepackage{bookmark}}{\usepackage{hyperref}}
\IfFileExists{xurl.sty}{\usepackage{xurl}}{} % add URL line breaks if available
\urlstyle{same} % disable monospaced font for URLs
\hypersetup{
  pdftitle={New Norwegian tangle and sugar kelp models},
  pdfauthor={Guri Sogn Andersen \& Kristina Kvile},
  hidelinks,
  pdfcreator={LaTeX via pandoc}}

\title{New Norwegian tangle and sugar kelp models}
\author{Guri Sogn Andersen \& Kristina Kvile}
\date{4 8 2022}

\begin{document}
\maketitle

\hypertarget{journal}{%
\section{Journal}\label{journal}}

Hva med Ecography? Diskuterte også med HGU å dele opp i to artikler. 1)
Artikkel der vi publiserer modellen fra blått karbon-prosjektet, og 2)
artikkel der vi har forenklet modell og viser fortid- og
framtidsprediksjoner. Det er nr 1 vi har skrevet outline til her:

\hypertarget{outline}{%
\section{Outline}\label{outline}}

Try to keep it short and simple. Avoid all the standard stuff about kelp
and the importance of kelp forests. Should go without explicit saying
when describing the Blue Carbon project aims etc. (The people seeking
this information probably knows all that stuff anyways \ldots)

\hypertarget{introduction}{%
\subsection{Introduction}\label{introduction}}

\begin{itemize}
\tightlist
\item
  Nordic Blue Carbon Project - background, aim, prospects?
\item
  Knowledge sought by Norwegian and Nordic authorities
\item
  The SDM models - uses and output
\item
  The kelp models

  \begin{itemize}
  \tightlist
  \item
    Unique in spatial coverage and resolution
  \item
    Potential uses
  \end{itemize}
\end{itemize}

\hypertarget{mat-met---translate-from-report-and-use-some-text-from-nordic-paper}{%
\subsection{Mat \& Met - translate from report and use some text from
Nordic
paper}\label{mat-met---translate-from-report-and-use-some-text-from-nordic-paper}}

\begin{itemize}
\tightlist
\item
  Observation data: Description of data, maps of observation points
\item
  Environmental data: Table with variable, source and summary stats
  (range, mean, etc \ldots)
\item
  Brief discussion of attempt of including urchin grazing pressure
\item
  Modeling: BRT fitting
\end{itemize}

\hypertarget{results-disc}{%
\subsection{Results + Disc}\label{results-disc}}

\begin{itemize}
\tightlist
\item
  Validation
\item
  Some key numbers and discussions of their validity. See Gundersen et
  al 2022 to avoid conflicts and/or duplicate publ
\item
  Strengths and weaknesses
\item
  Future prospects - ``teaser''
\end{itemize}

\hypertarget{authors}{%
\section{AUTHORS:}\label{authors}}

\textbf{Guri Sogn Andersen}, Kristina Kvile, \ldots., Helene Frigstad,
Hege Gundersen

\hypertarget{introduction-1}{%
\section{INTRODUCTION}\label{introduction-1}}

Kelp forests play an important role in the Norwegian carbon cycle, and
has been shown likely to sequester a significant amount of carbon each
year (Gundersen et al., 2011). Publications with global perspectives on
the important role of kelp in the marine carbon cycle \ldots.
(Krause-Jensen \& Duarte, 2016), recently sparked interest among Nordic
environmental authorities. This served as a backdrop for an initiative
from the Nordic Council of Ministers, seeking to update and increase the
knowledge base for the Nordic marine carbon cycle in order to include
blue forests in the national carbon inventories. This initiative birthed
The Nordic Blue Carbon Project, which aimed to create an updated,
quantitative overview of the carbon cycles of kelp, eelgrass and
rockweed in the Nordic region and address the knowledge gaps in
long-term carbon storage of macrovegetation in Nordic waters. A
fundamental part of this project was therefore to establish species
distribution models (SDMs) of the most important forest systems.
\emph{Kelp forests is by far the most abundant \ldots{} blabla ? }, and
the main results were SDMs of two kelp forest systems on both the Nordic
and Norwegian level. The Nordic model was recently published (Kvile et
al., 2022), and shows the spatial distribution of kelp forests
(presence/absence) in the Nordic region. The Norwegian data set mainly
stems from the national mapping of tangle kelp in the Norwegian Program
for Mapping Biodiversity (Bekkby et al., 2013), and includes information
on kelp density in an uniform way, which facilitated a much more
detailed modelling approach. To analyze and predict the distribution of
kelp based on a large number of potentially interacting environmental
variables, we used boosted regression trees (BRT), a method that
combines the advantages of machine learning and statistical regression
techniques (Elith et al., 2008). The here presented models show the
spatial variation in plant density within the two most important forest
forming kelp species, namely \emph{Saccharina latissima} (sugar kelp)
and \emph{Laminaria hyperborea} (tangle kelp), along the Norwegian
mainland. The prediction maps covers the entire coastline of Norway at a
25 m spatial resolution. These maps are therefore uniquely detailed, to
a level that is to our knowledge unprecedented in marine mapping.

\textbf{Skal vi ha med konvertering til biomasse?}

\hypertarget{materials-and-methods}{%
\section{MATERIALS AND METHODS}\label{materials-and-methods}}

\hypertarget{datasets}{%
\subsection{Datasets}\label{datasets}}

\hypertarget{kelp-registrations}{%
\subsubsection{Kelp registrations}\label{kelp-registrations}}

The data were sampled through the Norwegian Program for Mapping
Biodiversity (Bekkby et al., 2013). The program was designed for tangle
kelp (n = 11 891 observations) primarily, but sugar kelp was also
registered when found (n = 11 053). The dataset is thus well suited for
modelling tangle kelp, but less so in the case of sugar kelp. Data was
registered using an underwater video camera with an approximate
observation window of 5 square meters??? (\textbf{Se
kartleggingsprogram}). Density categories are quantified as absence,
single plants, scarce, moderately dense, and dense forest, representing
density categories of 0, 0.1, 0.5, 5 and 10 individual kelp plants per
m2 for tangle kelp and 0, 0.5, 1, 7 and 15 for sugar kelp (ref Bekkby?).

\emph{Print maps, remember contributions in caption (Leaflet option
turned off)}

\hypertarget{environmental-data-gis-models}{%
\subsubsection{Environmental data / GIS
models}\label{environmental-data-gis-models}}

In fitting the models for Norway, we used field-measured depth and a
range of environmental variables extracted from both national and global
spatial datasets.

Wave exposure was extracted from a model with a 25 m spatial resolution
(Isæus, 2004) based on fetch, wind speed and wind frequency. Ocean
current speed (m/s) was modeled at a horizontal resolution of 800 m
(NorKyst-800; Albretsen et al., 2011) using the three-dimensional
numerical ocean model ROMS (Shchepetkin \& McWilliams, 2005; Haidvogel
et al., 2008) and downscaled to 25 m resolution. The 90th percentiles
for seabed currents were used in our analyses. Due to the lack of a full
cover, high resolution substrate model, which does not exist for Norway,
slope (°) and terrain curvature (m) were derived from a bathymetry model
with a 25 m horizontal resolution, provided by the Norwegian
Hydrographic Service. Terrain curvature was estimated with a 250, 500
and 1 000 m spatial calculation window (Bekkby et al., 2009). Slope and
curvature were included as proxies for substrate under the assumption
that gentle slopes and depressions are areas of higher likelihood of
sedimentation and thus not hard substrate. Light at the seabed
(photosynthetically active radiation, PAR) has been modeled using data
from ocean color satellite sensor estimates (Saulquin et al., 2013,
Populus et al., 2017) at a resolution of 100 m, and was downscaled to a
model with 25 m resolution. All predictor variables given above,
including the 25 m bathymetry model, were available as full cover GIS
layers and used for making predictive distribution maps for kelp. We
also included variables from the global Bio-ORACLE datasets on
nutrients, dissolved oxygen, temperature and salinity at mean, maximum
and minimum depths (within grid cells of 5 arc xxx minutes resolution,
which corresponds to slightly less than 10 km at the equator). Historic
monthly data on seabed temperature and salinity at a grid resolution of
0.25° (\textbf{check}) was gathered from the Copernicus Marine Service,
and annual mean and max temperatures at each observation point was
calculated and used in fitting the model (\textbf{Check script or ask
Kristina}), while full-coverage Bio-ORACLE layers were used in the map
predictions of the present state of the density distributions.
\textbf{Something about how we coupled historic data in the fitting
process (from copernicus). Buffers etc \ldots{} We need perhaps to
justify the choice of min/mean or max depth in the BO-layers. For the
meantemp, the mindepth was used, while for the maxtemp the mean depth
was used. This was mainly based on the first version of the model, but
needs perhaps to be explained in some other way to simplify?} Even
though the resolution is relatively low, we chose to use Bio-ORACLE for
temperature and salinity data in the hope of utilizing their future
climate projection layers for these variables (corresponding to
predictions for the different RCP scenarios). It will eventually be
possible to plug these predictions into the model to make scenarios for
future kelp distribution, plant densities and biomass along the
Norwegian coast. \textbf{Compare resolution of Copernicusdata with
Bio-ORACLE. If we must justify the use of Bio-ORACLE in predictions,
maybe better to say something about coverage (which is far better with
BO).}

We coupled kelp observations (\emph{Figure X}) to environmental data by
extracting values from environmental raster layers cells in which the
kelp observations were located. If a kelp observation was outside the
cover of a raster layer, or data was missing in that specific cell, the
environmental variable was set to missing (NA).

\hypertarget{model-framework}{%
\subsection{Model framework}\label{model-framework}}

To predict the spatial distribution of kelp, boosted regression tree
(BRT) models were built using the gbm (Greenwell et al., 2020) and dismo
(Hijmans et al., 2020) libraries in R (R Core Team, 2020). When using
this approach, many simple models (``trees'') are built and combined
(``boosted'') for prediction (Elith et al., 2008). The BRT models have
several advantages, for example, a large number of predictor variables
of different types can be implemented without prior data transformation;
potential complex interactions between variables are automatically
handled; models are insensitive to outliers; and missing values are
accommodated (Elith et al., 2008). Moreover, BRTs can handle sharp
discontinuities, which is critical in cases such as ours, where the
species' distribution is small relative to the covered area.

Boosted Regression Tree models were built for Laminaria and Saccharina
separately, with kelp plant density as response variable. The response
was transformed (log + 1), wich supported the assumption of a gussian
error distribution. We kept most model parameters to their defaults
(Hijmans et al., 2020) but used an optimization routine based on
cross-validated deviance to determine the optimal combination of three
important parameters: learning rate (testing 0.0001, 0.001 and 0.01),
tree complexity (testing 2, 4 and 6), and bag fraction (testing 0.5, 0.6
and 0.7) (Smith, 2021). Learning rate sets the weight given to each
tree, and smaller learning rate is preferable but computationally costly
as a higher number of trees is required. Tree complexity defines the
complexity of interactions that can be fitted, if required by the data.
Bag fraction is the fraction of the data used in selecting variables at
each iteration (each tree), which makes the model stochastic and has
been shown to improve predictive power (Elith et al., 2008). The optimal
parameter combinations were learning rate 0.01, tree complexity 6 and
bag fraction 0.6 (7700 trees) for Laminaria, and learning rate X, tree
complexity X and bag fraction X (X trees) for Saccharina. After
identifying the optimal model with all explanatory variables included,
we ran a simplification procedure to remove variables that contributed
little to model predictive power using k-fold cross validation (Elith et
al., 2008).

\hypertarget{final-model}{%
\subsection{Final model}\label{final-model}}

Based on the final BRT models with unimportant predictor variables
removed, we predicted the density of kelp along the Norwegian coast
using the `predict' function in the raster package in R (Hijmans, 2022).
\textbf{Noe om konvertering til biomasse??}

We created a uniform raster object of all environmental predictor
variables included in the models, resampling all GIS-layers to a 25 m
resolution covering the whole model area (the Norwegian coast). While
measured depth from the kelp registration was used in the model fitting,
we used the full cover bathymetric data (m) for the predictions. And
while historic temperature data was used in the fitting process, a
static representation of present day annual mean and maximum temperature
based on long term averages (Bio-Oracle) was used for the spatial
predictions (\textbf{sjekk spesifisering av variabel}). Predictions from
BRT models are given for all cells regardless of the presence of missing
values for some variables, and missing values are set to the overall
mean value for that variable.

\hypertarget{results-and-discussion}{%
\section{Results and discussion}\label{results-and-discussion}}

We here present uniquely detailed models of the distribution of the two
most important forest forming kelp species in Norway, namely sugar kelp
(\emph{Saccharina latissima}) and tangle kelp (\emph{Laminaria
hyperborea}). The underlying kelp data predominantly stem from surveys
conducted to document the distribution of tangle kelp, and data from
typical sugar kelp areas are therefore underrepresented. This may have
lead to observations of sugar kelp being biased towards low density
registrations, which may contribute to an underestimation of sugar kelp
densities, but not necessarily of the total distribution area.

The tangle kelp model performed reasonably well, with a total deviance
explained as close to 75\% (a measure of variation in the data explained
by the model) and correlation values between data and model predictions
at 0.87 (based on training data) and 0.79 (based on cross-validation,
CV). The sugar kelp model performed less well. Although the total
deviance explained was approximately the same, the correlation values
were much lower (X and X for training data and CV, respectively)
\textbf{Maybe move these numbers to a table}

The most important explanatory variables for tangle kelp were bottom
depth and wave exposure contributing with 26.4 \% and 22 \% of the
explained deviance respectively (\textbf{check this}), followed by
nitrogen (11.8 \%), curvature (9.8 \%), dissolved oxygen (7.4 \%),
phosphate (5.2 \%), mean temperature (4.9 \%), max temperature (4.6 \%),
currents (4.3 \%) and finally max salinity (3.5 \%). All other variables
were cut in the simplification process. \emph{Add marginal effect plots
here} \emph{Perhaps also interaction plots?}

From report: When including all areas with densities ≥1 plant per m2, we
estimated the total area covered by tangle kelp to be 18 155 km2, and 44
300 km2 for sugar kelp. However, when including only high-density areas,
by setting cut-off values for forest at ≥5 for tangle kelp and ≥7 for
sugar kelp, these numbers were reduced to 3 810 km2 and 3 607 km2. These
numbers are slightly lower for tangle kelp and higher for sugar kelp
compared to former estimates by Gundersen et al.~(2011), who estimated 5
900 km2 and 2 000 km2 for tangle kelp and sugar kelp, respectively,
using a rule-based GIS model. It should also be taken into consideration
that, since Gundersen et al.~(2011) made these estimates, there has been
a significant regrowth of kelp (mostly sugar kelp) following sea urchin
destruction. The new kelp distribution models provided very detailed
maps of densities, which varied greatly depending on environmental
factors (Figure 2 and Figure 3). Insight into this spatial variation
showed that there are vast areas with low densities of kelp forest along
the Norwegian coast, compared to the high-density forests (Figure 4).
The primary production is calculated and presented in Chapter 3.

\end{document}
